\gddchapter{Interview sequence}


Below is the guideline for the semi-open guided interview. 
The interview starts with easy warm-up questions and progresses towards harder questions as it goes along. 

The answers were translated from polish back to english.

\section{Introduction}
Remind the participant that this is not a test, rather a joint exploration.

\section{Icebreakers}

\begin{QandA}
   \item Now that you’ve tried both versions, how are you feeling? Was there anything that immediately surprised you?
         \begin{answered}
               Marcel: Hmm, tired.
               My eyes are tired because my eyesight is so bad and I couldn't read.


               Gabriel: The letters are too small anyways.


               Marcel: It's not your fault here.


               I'm a little bit tired because I'm really not that much into story-like games, so it was a little bit tiring, but I suppose that's just my unique experience, I guess, in that sense.


               Gabriel: Yeah, I mean, it doesn't make it a bad experience, don't worry about it.
               And was there immediately something that, like, anything that, like, immediately surprised you when you were playing?


               Marcel: Surprised, I mean, let's think.
               I guess that's an answer because I can think of something right away.
            \end{answered}

   \item If you had to describe the difference between these two games to a person who hasn't played either, what would you say?
         \begin{answered}
               Okay, so the voice-operated game
               left me with much more imagination in that sense that I was being creative about my actions maybe just a little bit but still it made me feel more involved I feel like and the second game was different in that sense that
               I've had clear actions that I could do and I could navigate the field much faster but
               Yeah, I've had less use of my imagination in the second game.
               So I guess the main difference would be imagination in that sense.
               I was being less imaginative in the second game, like in the text input game.
               Not voice input, but keyboard input.
         \end{answered}
\end{QandA}




\section{Grand tour questions}
(remember to pause here 3-5 seconds)

\begin{QandA}


   \item From the moment you first used your voice to navigate, what were you thinking as you decided what to say to the system?

         \begin{answered}
            I was thinking about making an answer but in a fun way to somewhat test the system.
            That's an honest answer.
         \end{answered}

    \item How was it like to use the voice input in the game?
         \begin{answered}
            It was fine.
            It may be a little bit tiring after a while but in the time span that we had tested I didn't find that it was tiring but I was just slow with my inputs
            And that's okay, I was more immersed into the situation than the keyboard version
         \end{answered}


\end{QandA}




\section{Core comparative questions}

\begin{QandA}
   \item In which version did you feel more 'inside' the story world?
         \begin{answered}
            Marcel: Yeah, I definitely felt more immersed in the voice controlled version of the game.
            even if the system didn't necessarily included my little fun quirky additions to the inputs.


            Gabriel: I encourage you to be poetic about it.


            Marcel: Yes, even if it didn't include or even mention nothing about my ambitions, it was way more fun in that sense.
            I was way more immersed in the story.
            I really cared about my friend.
            I didn't leave no trace or anything, true, true, true, but with the keyboard version I just was moving, just doing robotic tasks and basically speedrunning the game.

           
         \end{answered}

   \item How did the 'effort' of thinking of what to say naturally compare to the 'effort' of knowing which keys to press?

         \begin{answered}
            Well, thinking what to say takes more effort than pressing a key, obviously.
            I mean, maybe it's not obvious, but...
            It takes more effort because you don't have your actions laid out in front of you and you have an open field, creative open field as of what you can really do.
            specific number of inputs or actions you can do and you can navigate it more quickly so it takes much less effort I think.
         \end{answered}

    \item Did the freedom to say 'anything' make you feel more in control, or was it more overwhelming than having a fixed list of keys?

         \begin{answered}
            Well...
            It gave me a sense of more freedom.
            But when I saw that my actions are fairly limited, I felt like my freedom was taken away from me a little bit.
            I don't think...
            I can see how it could be a little bit overwhelming.
            It was a little bit because I was moving slow and processing a lot of stuff and thinking about what I'm going to say.
            But I think the freedom was winning in that battle.
            Freedom opposed to being overwhelmed.
         \end{answered}
    \item How would you compare the overall experience of using the two game versions? Which one would you choose and why?

         \begin{answered}
            Marcel: So, overall experience...
             In the first game I've had, the voice activated game, I've had much more fun, to be honest.
             was being creative, I was being silly, I was thinking about my actions and that was way more engaging in that sense.
            In the keyboard input game I had my actions laid out in front of me so
            It was way more of a speedrunning contest for me, like I wanted to do things as quickly as I can.
            I sort of started to skip the storyline and just...
            I had my options in front of me, so I knew kind of what to do.
            And it became way less fun.
            So I had to create the fun myself in trying to do the game faster.
            And that was my engagement.
            So I think, like,
            If you, as a developer, are wanting to get the engagement of just the story and just the game, how you design it, the first one, in my experience, was way more fun and engaging and interesting because I had to think about it.
            Gabriel: So which one would you choose then?
            I would choose the first one.


            Gabriel: Okay.
            Because of?


            Marcel: Because of it being more interesting.
            more engaging because I don't use normally any sort of voice input utilities and I don't play games that include that sort of thing.
            So it was something new for me, something more fun, more interesting, not the same old, same old.
            compared to the keyboard type games.
            

            
         \end{answered}
\end{QandA}






\section{Probing}
In this part of the interview look into the sticky moments noted above and ask follow-up questions probing for more detail.
Include lexical reflection (eg. you used the word "clunky" when explaining navigation. Could you tell me more about what was happening
in that moment?)

% The probing was done as follow up questions to the pre-existing question list rather than separate questions in this test.

\begin{QandA}
   \item It looks like you really value creativity when it comes to playing games. Would you agree with that?
   \begin{answered}
      I would agree with that because
      If I played, let's say, a lot of games, then they can become kind of dull.
      So if something is new for me, or it's like, you know, creative or interesting,
      I value that because I'm pretty worn out by some normal or usual type of games like the second one.
      It was way more similar to the games that I played in the past, I would say.
   \end{answered}
\end{QandA}


\section{Final reflection}

\begin{QandA}
   \item After reflecting on everything today, is there anything else you’d like to add about using your voice to play this game that we haven't talked about?

         \begin{answered}
            Marcel: Well, I felt like I was a little bit limited by my eyesight and the small font was holding me back because most of my efforts was made to even read the text.
            Marcel: But that's the personal issue, I guess.


            Gabriel: No, no, the text is too small.  It's not just your issue.


            Marcel: I guess it's aesthetic.

            Gabriel: No, no, it's just my poor choice. Please carry on.


            Marcel: So if I were to make a change or add something, I would add
            just change the text to be bigger or more easily read.
            I was getting a little bit tired by reading so much text about like
            describing the scene yeah it was very specific a lot of things and I get that it wants to create a scene maybe I'm lazy in that sense that when I read so much of it I get a little bit overwhelmed so I don't know if that's something to add but this is like a comment I
            I don't think I have anything in mind specifically.
         \end{answered}
\end{QandA}

\section{Closing}
Thank the participant for his time and tell them again how much value his input and time gives me.






